(1) 圆柱坐标散度为 \(r^{-1}\partial_r(rv_r)+\partial_zv_z=-\alpha+\alpha=0\)。轴向涡量方程 \((\boldsymbol v\cdot\nabla)\omega=(\boldsymbol\omega\cdot\nabla)v_z+\nu\nabla^2\omega\) 在本构型化为
\[
-\frac{\alpha r}{2}\omega'=\alpha\omega+\nu(\omega''+\omega'/r).
\]
左边为径向平流，右边分别为涡伸长和黏性扩散。
(2) 设 \(\omega=C e^{-br^2}\)。代入上式并分别比较常数及 \(r^2\) 项得 \(b=\alpha/(4\nu)\)。由 \(\Gamma=2\pi C\int_0^\infty re^{-br^2}dr=\pi C/b\)，故
\[
\omega(r)=\frac{\Gamma\alpha}{4\pi\nu}\exp[-\alpha r^2/(4\nu)].
\]
该高斯同时使 \(r^0\) 和 \(r^2\) 系数为零，满足方程及有限环量。
(3) 从 \(\partial_r(rv_\theta)=r\omega\) 且轴心速度有限积分：
\[
v_\theta(r)=\frac{\Gamma}{2\pi r}(1-e^{-br^2}),\quad
\Gamma(r)=2\pi rv_\theta=\Gamma(1-e^{-br^2}).
\]
小 \(r\) 展开 \(v_\theta\simeq \Gamma br/(2\pi)=\Gamma\alpha r/(8\pi\nu)\)，为刚体旋转核；大 \(r\) 得 \(v_\theta\simeq\Gamma/(2\pi r)\)。
(4) 高斯积分
\[
\mathcal Z=2\pi C^2\int_0^\infty re^{-2br^2}dr=\frac{\pi C^2}{2b}=\frac{\Gamma^2\alpha}{8\pi\nu}.
\]
故涡旋贡献的量纲为 \([\mathcal Z]=L^2/T^2\)（每单位轴向长度的面积积分），\(\nu\to0^+\) 时涡核宽度 \(\sqrt{4\nu/\alpha}\to0\)，\(\mathcal Z\to\infty\) 而 \(\nu\mathcal Z\) 有限。外加线性拉伸场的应变在无限截面不衰减，其总黏性耗散本身发散；不能用局部涡量积分代替全域应变耗散。
